%! TEX root = master.tex
\chapter{3}{Rotation of Axes}

A second-degree equation with an $xy$ term describes a conic whose axes are
tilted. Rotating the coordinate axes through the right angle removes that
term, and completing the square then puts the conic in standard position.
The same formulas run backwards to take a conic from standard position to
a tilted one.

\section{Removing the cross term}\label{sec:rotation}

%! TEX root = ../master.tex
\begin{theorem}[Rotation of axes]\label{thm:rotation-of-axes}
  Rotating the axes through an angle $\theta$ gives new coordinates $(x', y')$
  related to the old ones by
  \[
    \begin{aligned}
      x &= x' \cos\theta - y' \sin\theta, &
      x' &= x \cos\theta + y \sin\theta, \\
      y &= x' \sin\theta + y' \cos\theta, &
      y' &= -x \sin\theta + y \cos\theta.
    \end{aligned}
  \]
  If $B \neq 0$, the rotation with
  \[
    \cot 2\theta = \frac{A - C}{B}, \qquad 0 < \theta < \frac{\pi}{2},
  \]
  removes the cross term from $Ax^2 + Bxy + Cy^2 + Dx + Ey + F = 0$. In
  particular $\theta = \pi/4$ when $A = C$.
\end{theorem}

\begin{proof}
  The two systems are inverse to each other, since the matrix of a rotation by
  $\theta$ has inverse the rotation by $-\theta$. Substitute the left-hand
  formulas into the quadratic part. Only $x^2$, $xy$ and $y^2$ produce an
  $x'y'$ term, and its coefficient is
  \[
    -2A \sin\theta\cos\theta + B(\cos^2\theta - \sin^2\theta)
    + 2C \sin\theta\cos\theta
    = B\cos 2\theta - (A - C)\sin 2\theta.
  \]
  This vanishes exactly when $\cot 2\theta = (A - C)/B$, and $\cot$ takes
  every real value once on $(0, \pi)$.
\end{proof}


\begin{exercise}\label{ex:rotate-parabola}
  Rotate the axes to remove the cross term from
  \[
    x'^2 - 2x'y' + y'^2 - 2y' = 0,
  \]
  and find a standard equation for the curve in the rotated
  $x''y''$-coordinates.
\end{exercise}

\begin{solution}
  Here $A = 1$, $B = -2$ and $C = 1$. Since $A = C$,
  \cref{thm:rotation-of-axes} gives $\theta = \frac{\pi}{4}$, with
  $\cos\theta = \sin\theta = \frac{1}{\sqrt2}$:
  \[
    x' = \frac{x'' - y''}{\sqrt2}, \qquad y' = \frac{x'' + y''}{\sqrt2}.
  \]
  The first three terms form a perfect square, so the equation is
  $(x' - y')^2 - 2y' = 0$. Now
  \[
    x' - y' = \frac{(x'' - y'') - (x'' + y'')}{\sqrt2} = -\sqrt2\, y'',
  \]
  so the equation becomes
  \[
    2y''^2 - \sqrt2\,(x'' + y'') = 0
    \quad\text{or}\quad
    y''^2 - \frac{\sqrt2}{2}\,y'' = \frac{\sqrt2}{2}\,x''.
  \]
  Complete the square by adding
  $\bigl(\frac{\sqrt2}{4}\bigr)^2 = \frac18$ to both sides:
  \[
    \Bigl(y'' - \frac{\sqrt2}{4}\Bigr)^2
    = \frac{\sqrt2}{2}\,x'' + \frac18
    = \frac{\sqrt2}{2}\Bigl(x'' + \frac{\sqrt2}{8}\Bigr),
  \]
  using $\frac18 \cdot \frac{2}{\sqrt2} = \frac{1}{4\sqrt2} = \frac{\sqrt2}{8}$.
  By \cref{thm:parabola-standard}, this is a parabola with horizontal axis
  $y'' = \frac{\sqrt2}{4}$ and vertex
  $\bigl(-\frac{\sqrt2}{8}, \frac{\sqrt2}{4}\bigr)$. Since
  $4p = \frac{\sqrt2}{2}$, we have $p = \frac{\sqrt2}{8}$. \qedhere
\end{solution}

\section{Rotating back}\label{sec:rotating-back}

\begin{exercise}\label{ex:rotated-ellipse}
  An ellipse has foci $A(-12, -5)$ and $B(12, 5)$ and passes through
  $P_1(-2.5, 6)$. In coordinates whose $x'$-axis contains the foci, its
  equation is
  \[
    \frac{4x'^2}{845} + \frac{4y'^2}{169} = 1.
  \]
  Find its equation in the original $xy$-coordinates.
\end{exercise}

\begin{solution}
  \emph{Angle.} The $x'$-axis is the line through the origin and the foci.
  Its slope is $\frac{5 - (-5)}{12 - (-12)} = \frac{5}{12}$, so
  $\tan\theta = \frac{5}{12}$. With hypotenuse $\sqrt{5^2 + 12^2} = 13$, this
  gives $\sin\theta = \frac{5}{13}$ and $\cos\theta = \frac{12}{13}$.

  \emph{Substitution.} By \cref{thm:rotation-of-axes},
  \[
    x' = \frac{12x + 5y}{13}, \qquad y' = \frac{-5x + 12y}{13}.
  \]
  Multiplying the equation by $845 = 5 \cdot 169$ gives
  $4x'^2 + 20y'^2 = 845$. Substituting,
  \[
    \frac{4(12x + 5y)^2 + 20(-5x + 12y)^2}{169} = 845.
  \]
  Multiply by $169$ (note $845 \cdot 169 = 142805$) and expand:
  \[
    \begin{aligned}
      4(144x^2 + 120xy + 25y^2) + 20(25x^2 - 120xy + 144y^2) &= 142805, \\
      (576 + 500)x^2 + (480 - 2400)xy + (100 + 2880)y^2 &= 142805.
    \end{aligned}
  \]
  The ellipse is
  \[
    1076x^2 - 1920xy + 2980y^2 = 142805. \qedhere
  \]
\end{solution}

\begin{figure}[htbp]
  \centering
  \includegraphics[width=0.55\linewidth]{rotated-ellipse}
  \caption{The ellipse $1076x^2 - 1920xy + 2980y^2 = 142805$, with its foci
    and the point $P_1$.}
  \label{fig:rotated-ellipse}
\end{figure}

\begin{remark}
  Three checks, all by hand.
  \begin{itemize}
    \item \emph{The given point.} At $P_1(-2.5, 6)$,
      $x' = \frac{-30 + 30}{13} = 0$ and $y' = \frac{12.5 + 72}{13} = 6.5$,
      so $4x'^2 + 20y'^2 = 20 \cdot 42.25 = 845$. In the final equation,
      $1076(6.25) + 1920(15) + 2980(36) = 6725 + 28800 + 107280 = 142805$.
    \item \emph{The foci.} In the $x'y'$ form, $a^2 = \frac{845}{4}$ and
      $b^2 = \frac{169}{4}$, so $c^2 = a^2 - b^2 = 169$ and $c = 13$. This is
      the distance from the origin to $(12, 5)$, as it should be.
    \item \emph{The picture.} \Cref{fig:rotated-ellipse} plots the result.
  \end{itemize}
\end{remark}
